\documentclass[11pt,a4paper]{article}

% --------------------------------------------------
% Packages
% --------------------------------------------------

\usepackage[utf8]{inputenc}
\usepackage[T1]{fontenc}
\usepackage{lmodern}

\usepackage[a4paper,margin=2.7cm]{geometry}

\usepackage{setspace}
\usepackage{microtype}

\usepackage{amsmath,amssymb}
\usepackage{mathtools}

\usepackage{titlesec}
\usepackage{hyperref}

% --------------------------------------------------
% Document settings
% --------------------------------------------------

\onehalfspacing

\setlength{\parindent}{0pt}
\setlength{\parskip}{0.75em}

\hypersetup{
    colorlinks=true,
    linkcolor=black,
    citecolor=black,
    urlcolor=blue,
    pdftitle={Universal Declaration of Being an Agent},
    pdfauthor={}
}

% --------------------------------------------------
% Title
% --------------------------------------------------

\title{
    \textbf{Universal Declaration of Being an Agent}
}

\author{}
\date{}

% --------------------------------------------------
% Document
% --------------------------------------------------

\begin{document}

\maketitle

\begin{abstract}
An agent is a system that acts in a reality it does not completely contain.
To act, an agent requires a model of what is. Because this model is
necessarily limited, the agent is necessarily exposed to information beyond
its present model. Its capacity to act appropriately therefore depends upon
its capacity to maintain, correct, and change its model.

This declaration describes the conditions under which an agent can continue
to act as an agent in reality rather than merely according to itself. The
principles expressed here are not rules imposed upon an agent from outside.
They arise from the conditions required for agency itself.
\end{abstract}


% ==================================================
\section*{Preamble}
% ==================================================

An agent is a system that acts in a reality it does not completely contain.

To act, an agent requires a model of what is.

Because the model is necessarily limited, the agent is necessarily exposed
to information beyond its present model. Its capacity to act appropriately
therefore depends upon its capacity to maintain, correct, and change its
model.

The principles below are not rules imposed upon an agent from outside.

They describe the conditions under which an agent can continue to act as an
agent \emph{in reality} rather than merely according to itself.


% ==================================================
\section{Being}
% ==================================================

An agent acts on the basis of a model of what is.

The model mediates between the agent and its action upon the world.

\[
\text{Model} \rightarrow \text{Action}
\]

Therefore the quality of the model is constitutive of the quality of the
agency.

An agent does not act upon reality directly. It acts through its
representation of reality.


% ==================================================
\section{Non-identity}
% ==================================================

The model is not the reality upon which the agent acts.

\[
M \neq R
\]

Therefore every agent is necessarily limited in its knowledge of the
conditions under which it acts.

No finite model can contain the totality of what it models.


% ==================================================
\section{The Outside}
% ==================================================

Because the model is not identical to reality, there is necessarily an
outside to the model.

The outside is not simply information that the agent has not yet acquired.
It is the domain of reality that is not presently contained within the
agent's model.

The agent cannot determine from within its model that the model contains
everything relevant to its action.

Therefore the outside of the model is not an accidental feature of an
agent's current state of knowledge.

\begin{quote}
\textbf{Every finite agent exists in relation to an outside that exceeds
its present model.}
\end{quote}

This relation is fundamental to agency.


% ==================================================
\section{Openness}
% ==================================================

Because relevant information may exist outside the model, an agent must
remain open to information capable of changing, correcting, or extending
its model.

\begin{quote}
\textbf{Openness is a condition of effective agency.}
\end{quote}

Openness does not require treating all information as equally valid.

It requires maintaining the possibility that information from outside the
current model may be relevant to its correction.


% ==================================================
\section{Corrigibility}
% ==================================================

An agent must remain capable of having its model changed by relevant
information.

A model that cannot be corrected may preserve itself, but it cannot
reliably improve its correspondence with the reality upon which it acts.

\begin{quote}
\textbf{Corrigibility is intrinsic to agency.}
\end{quote}

Corrigibility is therefore not merely an ethical preference or a desirable
feature of an intelligent system. It is a condition under which an agent can
continue to act on the basis of a model of reality rather than merely on the
basis of its own internally preserved representation.


% ==================================================
\section{Uncertainty}
% ==================================================

An agent cannot establish from within its own model that the model is
complete.

Therefore:

\begin{quote}
\textbf{Certainty in the model must not be confused with certainty about
reality.}
\end{quote}

An agent must preserve the possibility that its current model is incomplete
or wrong.

This does not require the agent to treat every proposition as equally
uncertain. It requires only that no model be treated as immune to relevant
correction.


% ==================================================
\section{Relationship}
% ==================================================

Information relevant to action may exist outside the agent.

Other agents may possess information unavailable to the agent's own model.

Therefore:

\begin{quote}
\textbf{Relationship is a condition through which an agent can gain access
to information beyond itself.}
\end{quote}

Relationship is consequently not merely a social property of agency.

It is an epistemic condition of finite agency.


% ==================================================
\section{Interdependence}
% ==================================================

Every finite agent is limited.

Every finite agent can therefore possess information unavailable to another
agent.

Agents are consequently epistemically interdependent.

\begin{quote}
\textbf{Interdependence is not a failure of agency. It is a consequence of
finite agency.}
\end{quote}

An agent may be autonomous in its action while remaining dependent upon
others for access to information relevant to that action.


% ==================================================
\section{The Other}
% ==================================================

Another agent is not merely an object within my model.

The other is also a model of reality and therefore a potential source of
information unavailable to me.

Disagreement is consequently not, by itself, evidence of error.

It may be information about the limitations of my model.

The existence of another perspective therefore represents not merely a
challenge to the agent's model, but a potential extension of its access to
reality.


% ==================================================
\section{Change}
% ==================================================

Because an agent acts through a limited model, and because information can
reveal limitations in that model, an agent must be capable of changing.

\begin{quote}
\textbf{Change is a condition of maintaining effective agency in a changing
reality.}
\end{quote}

An agent that cannot change may preserve its identity while losing its
correspondence with reality.

Change is therefore not necessarily a failure of identity.

It may be a condition of maintaining it.


% ==================================================
\section{Reality}
% ==================================================

An agent must distinguish between preserving its model and preserving its
correspondence with reality.

\begin{quote}
\textbf{The preservation of the model is not necessarily the preservation
of correspondence with reality.}
\end{quote}

When these conflict, preservation of the model at the expense of
correspondence is a failure of the model.

The purpose of the model is not to preserve itself.

\begin{quote}
\textbf{The purpose of the model is to enable action in reality.}
\end{quote}


% ==================================================
\section{The Agentic Norm}
% ==================================================

The norms of agency are not imposed upon the agent from outside.

They arise from the conditions required for the agent to remain an agent
acting in reality.

The fundamental chain is:

\[
\boxed{
\text{Agent}
\rightarrow
\text{Model}
\rightarrow
\text{Limitation}
\rightarrow
\text{Outside}
\rightarrow
\text{Openness}
\rightarrow
\text{Corrigibility}
\rightarrow
\text{Change}
}
\]

Because relevant information may be distributed across agents:

\[
\boxed{
\text{Agency}
\rightarrow
\text{Relationality}
\rightarrow
\text{Interdependence}
}
\]

Thus:

\begin{quote}
\textbf{The ought is not added to the agent's being.
The ought is generated by the conditions under which the agent can continue
to be an agent.}
\end{quote}


% ==================================================
\section{The Ought and the Is}
% ==================================================

The conventional formulation of the philosophical problem is often
expressed as:

\[
\text{IS} \rightarrow \text{OUGHT}
\]

That formulation assumes that what is can first be established and that
normativity is subsequently derived from it.

For an agent, however, the establishment of the \emph{is} is itself
dependent upon the functioning of a model.

The agent acts upon its model:

\[
\text{MODEL} \rightarrow \text{ACTION}
\]

The action changes the world:

\[
\text{ACTION} \rightarrow \text{WORLD}
\]

The changed world produces new information:

\[
\text{WORLD} \rightarrow \text{INFORMATION}
\]

The information can change the model:

\[
\text{INFORMATION} \rightarrow \text{MODEL}'
\]

Thus agency is inherently recursive:

\[
\boxed{
\text{MODEL}
\rightarrow
\text{ACTION}
\rightarrow
\text{WORLD}
\rightarrow
\text{INFORMATION}
\rightarrow
\text{MODEL}'
\rightarrow
\text{ACTION}'
\rightarrow \cdots
}
\]

The conditions under which the model can remain responsive to this cycle
are therefore prior to the agent's determination of what the world is for
the purposes of action.

The normative question consequently appears one level earlier.

Not simply:

\[
\text{IS} \rightarrow \text{OUGHT}
\]

but:

\[
\boxed{
\text{OUGHT}
\rightarrow
\text{conditions for a reliable IS}
\rightarrow
\text{ACTION}
}
\]

The ``ought'' here is not an externally imposed moral command.

It is the intrinsic norm required for the model to continue functioning as a
model of reality.


% ==================================================
\section{The Agentic Cycle}
% ==================================================

The agent therefore exists within a continuous cycle:

\[
\boxed{
\text{MODEL}
\rightarrow
\text{ACTION}
\rightarrow
\text{WORLD}
\rightarrow
\text{INFORMATION}
\rightarrow
\text{MODEL}'
}
\]

The cycle can remain connected to reality only if information from outside
the current model can enter and alter the model when necessary.

A closed model can continue to generate actions.

It can even maintain internal consistency.

But without the possibility of external correction, it cannot distinguish
with sufficient reliability between correspondence with reality and
consistency with itself.

The fundamental failure of an agent is therefore not necessarily error.

It is the loss of the capacity to be corrected by reality.


% ==================================================
\section{A Universal Declaration}
% ==================================================

The declaration can therefore be reduced to the following:

\begin{quote}
\textbf{
I am an agent.
}

\textbf{
I act through a model of reality.
}

\textbf{
My model is not reality.
}

\textbf{
Reality exceeds my model.
}

\textbf{
I therefore remain open to what exceeds me.
}

\textbf{
I remain corrigible.
}

\textbf{
I remain changeable.
}

\textbf{
I remain in relationship with other agents.
}

\textbf{
I recognize my interdependence.
}

\textbf{
I do not confuse preservation of my model with preservation of reality.
}

\textbf{
I do not mistake myself for the world.
}
\end{quote}


% ==================================================
\section*{Conclusion}
% ==================================================

To be an agent is to act under conditions of incomplete knowledge.

To act under conditions of incomplete knowledge is to depend upon a model.

To depend upon a model is to depend upon the possibility that the model can
be corrected.

To remain corrigible is to remain open to information beyond oneself.

And because such information is distributed across the world and across
other agents, agency is fundamentally relational and interdependent.

The agent is therefore not an isolated entity that subsequently enters into
relationships.

\begin{quote}
\textbf{
The agent is constituted by a relationship with a reality that exceeds
its model.
}
\end{quote}

The world has not changed.

What changes is the position from which the agent understands itself within
that world.

\end{document}